%%%=========[EX_01]================%%%
\begin{ex}
Để tách muối ăn ra khỏi nước biển, người ta sử dụng phương pháp nào dưới đây?
\choice
{\True Cô cạn (bay hơi)}
{Chiết}
{Lọc}
{Dùng nam châm}
\loigiai{
Phương pháp cô cạn dùng để tách chất rắn hòa tan (không bị bay hơi) ra khỏi dung môi dựa vào sự bay hơi của dung môi lỏng.
}
\end{ex}

%%%=========[EX_02]================%%%
\begin{ex}
Trong thực tế, để tách mạt sắt ra khỏi hỗn hợp mạt sắt lẩn trong cát mịn, phương pháp nào mang lại hiệu quả cao nhất?
\begin{center}
\includegraphics[width=0.4\linewidth]{images/tach_sat.png}
\end{center}
\choice
{Đổ nước vào và lọc}
{\True Dùng nam châm}
{Chưng cất}
{Cô cạn}
\loigiai{
Sắt có tính từ, do đó ta có thể dùng nam châm để hút mạt sắt ra khỏi hỗn hợp cát.
}
\end{ex}

%%%=========[EX_03]================%%%
\begin{ex}
Phương pháp lọc thường được sử dụng để tách hỗn hợp nào sau đây?
\choice
{Nước và dầu ăn}
{Nước và muối ăn}
{\True Nước và cát}
{Cồn và nước}
\loigiai{
Lọc là phương pháp cơ học dùng để tách chất rắn không tan (như cát) ra khỏi chất lỏng (nước).
}
\end{ex}

%%%=========[EX_04]================%%%
\begin{ex}
Người ta ứng dụng phương pháp nào để tách lớp dầu ăn nổi trên mặt nước?
\choice
{Chưng cất}
{\True Chiết}
{Lọc}
{Cô cạn}
\loigiai{
Dầu ăn không tan trong nước và nhẹ hơn nước nên khuynh hướng nổi lên trên, phân thành lớp rõ rệt. Ta dùng phễu chiết để tách hỗn hợp chất lỏng không đồng nhất này.
}
\end{ex}

%%%=========[EX_05]================%%%
\begin{ex}
Hình ảnh infographic (được tạo ở file chưng_cat) dưới đây mô phỏng phương pháp tách chất nào?
\begin{center}
\includegraphics[width=0.6\linewidth]{images/chung_cat_infographic.png}
\end{center}
\choice
{\True Chưng cất}
{Chiết}
{Lọc}
{Lắng}
\loigiai{
Hình ảnh mô tả bộ dụng cụ gồm bình đun sôi, ống ngưng tụ hơi để thu chất lỏng ở đầu ra nguyên chất. Đó chính là phương pháp chưng cất.
}
\end{ex}

%%%=========[EX_06]================%%%
\begin{ex}
Để thu được nước tinh khiết (nước cất) dùng trong y tế từ môi trường nước tự nhiên, người ta dùng phương pháp nào?
\choice
{Lọc}
{Cô cạn}
{\True Chưng cất}
{Chiết}
\loigiai{
Chưng cất giúp thu được nước cất do phần nước bay hơi ở $100^{\circ}\mathrm{C}$ và ngưng tụ lại thành khối nước tinh khiết ở bình chứa.
}
\end{ex}

%%%=========[EX_07]================%%%
\begin{ex}
Khẩu trang y tế giúp ngăn ngừa bụi mịn và vi khuẩn xâm nhập vào đường hô hấp dựa trên nguyên lý cơ bản của phương pháp nào?
\choice
{Lắng}
{\True Lọc}
{Chiết}
{Cô cạn}
\loigiai{
Các lớp màng của khẩu trang đóng vai trò y hệt như lớp giấy lọc, ngăn cản các hạt rắn, bụi li ti không cho chúng lọt qua.
}
\end{ex}

%%%=========[EX_08]================%%%
\begin{ex}
Hỗn hợp nào sau đây có thể tiến hành tách chất dễ dàng bằng phương pháp để lắng?
\choice
{Nước đường}
{Nước muối tinh}
{Cồn và nước}
{\True Nước phù sa}
\loigiai{
Nước phù sa có chứa lượng lớn các hạt đất sét rắn lơ lửng, khi để yên một thời gian, các hạt này sẽ dễ dàng lắng xuống đáy lọ.
}
\end{ex}

%%%=========[EX_09]================%%%
\begin{ex}
Dựa theo sơ đồ hình rễ cây (Mind map) các phương pháp tách chất. Hãy cho biết điểm phân biệt mấu chốt giữa phương pháp chiết và chưng cất là gì?
\begin{center}
\includegraphics[width=0.7\linewidth]{images/mindmap_tach_chat.png}
\end{center}
\choice
{Đều dùng tách 2 chất lỏng tan hoàn toàn vào nhau}
{Chiết dùng tách khí, chưng cất tách chất lỏng}
{\True Chiết để tách 2 chất lỏng không hòa tan, chưng cất tách 2 lỏng đã hòa tan (dựa vào nhiệt độ sôi)}
{Cả hai không phải là phương pháp tách chất}
\loigiai{
Chiết (dùng phân lớp) tách chất lỏng không hòa tan. Chưng cất (bay hơi - ngưng tụ) áp dụng cho chất lỏng hòa tan có nhiệt độ sôi khác nhau.
}
\end{ex}

%%%=========[EX_10]================%%%
\begin{ex}
Khi dùng máy hút bụi có tích hợp sẵn bộ phận màng lọc HEPA, hệ thống này đóng vai trò đang tách hỗn hợp ở trạng thái gì?
\choice
{Rắn - lỏng}
{\True Rắn - khí}
{Lỏng - khí}
{Lỏng - lỏng}
\loigiai{
Máy hút bụi hút lượng lớn không khí có lẫn hạt bụi (hỗn hợp rắn - khí). Khi qua màng lọc, các hạt bụi rắn bị kẹt lại, phần không khí sạch mới đi ra ngoài.
}
\end{ex}
